\documentclass[11pt]{article}
\usepackage[margin=1.1in]{geometry}
\usepackage{amsmath,amssymb,amsthm,mathtools}
\usepackage[T1]{fontenc}
\usepackage{microtype}
\usepackage{booktabs}
\usepackage{graphicx}
\usepackage{algorithm}
\usepackage{algpseudocode}
\usepackage{enumitem}
\usepackage[colorlinks=true,linkcolor=black,citecolor=black,urlcolor=black]{hyperref}

\newcommand{\R}{\mathbb R}
\newcommand{\E}{\mathbb E}
\newcommand{\Q}{\mathbb Q}
\newcommand{\PP}{\mathbb P}
\newcommand{\dd}{\mathrm d}
\newcommand{\dlv}{\ell}
\newcommand{\abs}[1]{\lvert #1\rvert}
\newcommand{\ip}[2]{\langle #1,#2\rangle}
\newcommand{\one}{\mathbf 1}
\newcommand{\Var}{\operatorname{Var}}
\newcommand{\Gh}{\widehat G}
\newcommand{\Nh}{\widehat N}
\newcommand{\Rt}{\widetilde R}
\newcommand{\HB}[1]{{\color{blue}{#1}}}

\theoremstyle{plain}
\newtheorem{theorem}{Theorem}[section]
\newtheorem{proposition}[theorem]{Proposition}
\newtheorem{lemma}[theorem]{Lemma}
\newtheorem{corollary}[theorem]{Corollary}
\theoremstyle{definition}
\newtheorem{assumption}[theorem]{Assumption}
\theoremstyle{remark}
\newtheorem{remark}[theorem]{Remark}

\title{
    \HB{
    Stochastic Discrete  Local Volatility}
}

\author{}
\date{}

\begin{document}
\maketitle

\begin{abstract}
    \HB{
    We present a fully arbitrage-free option market model using the SANOS~\cite{SANOS2026} framework
    based in ideas of~\cite{Wissel2007}.

    The model is the product of a risk-neutral discrete state martingale defined on a fixed grid with stochastic ``discrete local volatilities'',
    and a simple SANOS-type Black~\&~Scholes smoothing process which allows smooth interpolation
    of option prices between grid points.

    We train the model on S\&P 500 Index options data.}
\end{abstract}

\begin{abstract}
\noindent
We study option market models on a fixed lattice of cash strikes and expiry
dates whose state is the vector of gridded risk-neutral densities together with
the discrete local variances. Positivity of the local variances is equivalent to
absence of static arbitrage including calendar spreads, so the admissible set is
an open orthant and $\lambda=\log\dlv$ ranges over a Euclidean space: a diffusion
on the code-book is statically arbitrage-free by construction.

Two results follow from the choice of chart. First, the spot is a read-off,
$S_t=x^{\!\top}q^j_t$ at every expiry, an identity of the chart holding at all
$t$ with no limit taken. At the front expiry we exhibit a family of front
volatilities under which the density collapses onto a single node, so the spot
lands on a strike and every gridded contract settles at its exact payoff. That
pinning is a property of the chosen volatility and not of the chart, which
represents every terminal spot in the grid range; in \cite{Wissel2007} it
follows from the chart whatever the coefficients. Second, the model has a unique global
strong solution on $[0,T_1]$, up to and including the first expiry, for every
bounded and locally Lipschitz volatility structure and any number of driving
factors, with no damping, no frozen wing, and no design rule. The only
requirement is that the front volatility be integrable in time, which the
natural collapse rate $(T_1-t)^{-1/2}$ satisfies. The quantity to be
controlled is $\eta=\psi/q$, the density volatility relative to the density; it
is bounded by the grid alone. Wissel \cite{Wissel2007} must assume a bound of
this kind; here it is a theorem.

The no-arbitrage drift restriction itself is not new. In the continuum it is
\cite{CarmonaNadtochiy2008}, and on this lattice it is \cite{Wissel2007}, whose
local implied variance coincides with the discrete local variance used here as
an algebraic identity on the grid. What the density chart adds is that the drift
becomes pointwise, $\alpha^j_i=-\sum_f\eta^{j,f}_i\beta^{j,f}_i$, with no matrix
inversion and no rank condition, the apparent over-determination being resolved
by conservation of mass and of the spot.
\end{abstract}

\section{Introduction}\label{sec:intro}

\HB{
This article presents an arbitrage-free market model for for smooth surfaces of European option prices.

The core architecture follows the SANOS approach~\cite{SANOS2026} which models the stock price as the product
of a time-and-state discrete process~$X$ and a smoothing process~$Y$, in our case Black~\&~Scholes.
The former is used to encode the option market using a discrete form of local volatilities as code book. That does not mean
that the model is a local volatility model -- we merely use the (discrete) local volatilities as a convenient representation of the option surface.
Compared to using implied volatilities the advantage is that mere positivity of the local volatilities is sufficient to ensure
the the option market is free of static arbitrage at all times. We have shown in~\cite{SANOS2026} that this representation
can be fitted to real life market data very efficiently using linear programming.

In~\cite{DYSANOS2026} we have shown how to estimate the dynamics such (discrete) local volatilities of such a model under the
statistical measure~$\mathbb{P}$ which fit very well to real-life observed option bid/asks.

The key
contribution in this paper is showing how to construct a corresponding risk-neutral measure~$\mathbb{Q}$ or, equivalently,
the risk-neutral drift of the model. Our approach is closely related to the seminal but unpublished work~\cite{Wissel2007}.
}

An option market model reverses the usual direction of construction: instead of
specifying a process for the underlying and deriving option prices from it, one
models the code-book of prices directly, so that a calibration is a change of
coordinates rather than an inverse problem. Two conditions must then be
maintained at once. The state must remain free of \emph{static} arbitrage, and
its dynamics must admit one equivalent local martingale measure common to the
underlying and to every traded option. The second is the option analogue of an
HJM drift restriction: the volatility structure may be specified freely, and the
drift is then determined. The lattice here is \emph{finite}, finitely many
strikes and finitely many expiry dates, which is what is traded.

\paragraph{What is already known.} Schweizer and Wissel treat the problem in the
continuum in two companion papers, one in each direction of the surface:
\cite{Schweizer2007} varies maturity at a single strike, and
\cite{SchweizerWissel2008} varies strike at a single maturity. In each the
code-book is chosen so that its positivity encodes absence of arbitrage, the
drift restriction is derived, and existence is proved. The multi-strike
multi-maturity case, which is ours, is stated as open in
\cite{SchweizerWissel2008}, with the expectation that a new parametrisation will
be needed. Carmona and Nadtochiy
\cite{CarmonaNadtochiy2008} give a continuum local-volatility code-book with a
drift restriction, but no existence theorem. On a finite lattice, and answering
the multi-maturity question, stands Wissel \cite{Wissel2007}; his construction is
the one this paper is closest to, and it is described separately below.

\paragraph{Modern approaches.} The question was reopened from the data side, by
the scarcity of option-market histories for training and evaluating hedging
strategies \cite{BuehlerEtAl2019DeepHedging}. Cohen, Reisinger and Wang
\cite{CohenReisingerWang2021} build neural stochastic differential equations on a
floating lattice and penalise a least-squares residual in place of the drift
restriction.
\cite{WieseEtAl2019OptionMarkets} simulate an equity option book with a
generative adversarial network whose state is the discrete local variance grid
compressed by principal components, so that static arbitrage is absent by
positivity of that grid; that is the code-book choice this paper inherits, and
\cite{WieseEtAl2021MultiAsset} extend it to a multi-asset universe.
\cite{DYSANOS2026} decode statically arbitrage-free surfaces from a
low-dimensional code, fit a linear state process to it, and test for dynamic
arbitrage after the fact, reporting it in $12.5\%$ of states at zero transaction
cost.

None of the four enforces the drift restriction exactly, and that is where this paper starts:
the cross-sectional constraints govern each date's surface and its terminal
payoff cone, and leave the transition of a fixed contract between two dates
unconstrained. The choice of grid is part of the same point. All four index
their nodes relative to the spot rather than in fixed cash terms, and on such a
grid a node is not one contract through time, so there is no single instrument
whose price is required to be a martingale and nothing for an HJM-type
restriction to constrain. Fixing the strikes in cash is what buys the
restriction here.

\paragraph{The closest prior work.} Wissel \cite{Wissel2007} works on a finite lattice of fixed cash strikes and
fixed expiry dates, parametrises it by a price level together with positive local
implied volatilities, derives the no-arbitrage drift restrictions, and proves
constructive strong existence and uniqueness for a class of coefficients. Behind
the front expiry his local implied variance is \emph{exactly} the backward
discrete local variance used here, an algebraic identity on the grid, and his
drift restriction is, after a change
of measure, exactly the one we state in Theorem~\ref{thm:drift}. We claim no
priority for it; see Appendix~\ref{app:wissel}, which gives the identification in
full. That working paper appears never to have been published in archival
form.

\subsection*{What this paper does}

We keep the same lattice and change the chart: the state is the vector of
gridded \emph{densities} together with the discrete local variances, rather than
prices together with quantiles of digitals. Three consequences follow, and they
are the contributions.

\begin{enumerate}[leftmargin=1.8em,itemsep=0.45em]
\item \textbf{The spot is a read-off, and it settles exactly.}
The spot is $S_t=K^{\!\top}q^j_t$ at every expiry $j$ (Proposition~\ref{prop:spot}),
an identity of the chart holding at all $t$. At the front
expiry Proposition~\ref{prop:collapse} exhibits a family of front volatilities
under which $q^1$ collapses onto a single node, so the spot lands on a strike
and every gridded contract settles at its exact payoff
(Section~\ref{sec:nonpinning}). That pinning is a property of the volatility
chosen and not of the chart: by Remark~\ref{rem:bary} the terminal states
the chart admits carry every spot in the grid range, whereas in
\cite{Wissel2007} the pinning follows from the chart whatever the coefficients
(Remark~\ref{rem:pinning}). We do not construct dynamics reaching an off-lattice
terminal spot; that is left as future work.

\item \textbf{Existence, with no condition on the wings.} For every bounded and
locally Lipschitz volatility structure $b$ and any number of driving factors, the
model has a unique strong solution on $[0,T]$ for every $T<T_1$, positive by
construction; it extends to $T_1$ (Remark~\ref{rem:horizon}) and, restarting at each
expiry with the settling front, to the whole ladder $[0,T_M]$, unique given that
front (Corollary~\ref{cor:patch}). The quantity that must be controlled is
$\eta=\psi/q$, the density volatility relative to the density, and its growth across
expiries is bounded by the grid alone: Theorem~\ref{thm:wing} gives $\Theta_j\le N+1$
uniformly in the state and in the volatility structure. \emph{Wissel \cite{Wissel2007} must assume a bound of this
kind; here it is a theorem.} The front expiry costs one hypothesis, and it is a mild one: $\eta^1$
must be integrable in $t$, not bounded, because $\eta^j$ is affine in $\eta^1$
and an It\^o equation needs an integrable drift rather than a bounded one. The
front volatility used here diverges at rate $(T_1-t)^{-p}$, and the two demands
made of $p$ leave a window: $p<1$ is what Assumption~\ref{ass:F} asks, while
$p\ge\tfrac12$ is what makes the front density collapse, so that every gridded
contract settles at its exact payoff (Proposition~\ref{prop:collapse}).

\item \textbf{The drift is pointwise, and the count works out.} In density
coordinates the drift equation is $N+2$ conditions on $N$ unknowns, and is
consistent because the left null space of $\widehat G_j$ is exactly
$\mathrm{span}\{\one,K\}$, conservation of mass and of the spot
(Lemma~\ref{lem:null}). The state-independent factor then cancels from both
sides and the drift is available nodewise,
$\alpha^{j,i}=-\ip{\eta^{j,i}}{\beta^{j,i}}$, with no matrix inversion and no
rank condition.
\end{enumerate}

\subsection*{Scope, and what is not claimed}

The over-determination of the third contribution is not intrinsic: it is created
by the choice of density coordinates and removed by the same choice, and in
price coordinates the system is square and the question does not arise.

The existence statement is for a fixed finite grid. The constant is
$(M-1)(N+1)$ and depends on the grid. The one
hypothesis, Assumption~\ref{ass:F}, is a condition on the front volatility and
not on the state, and it is sufficient rather than necessary.

\section{The code-book}\label{sec:setup}

We assume deterministic forwards and discount factors, which reduces to a driftless spot
with unit discount factors, cf.~\cite{BuehlerVolAndDivs}. The spot at $t=0$ is $S_0$.

The approach we are presenting relies on a combination of a time-and-state 
backbone process~$X$ and an additional smoothing process $Y$, following the~\cite{SANOS2026} approach to smoothing. However, contrary
to that model here we present a full stochastic local volatility model. The weight of the model is carried
by the former and echos ideas of~\cite{Wissel2007}.

We will define the backbone process on a grid of fixed expiries $0=T_0< T_1<\dots<T_M$ and strikes $0< K^0< K^1<\dots<S_0<\dots<K^N<K^{N+1}$.
We will use two time indices: $t$ will denote model time and $j$ will index fixed expires. We will use $i$ to index strikes
Contrary to~\cite{SANOS2026} and similar to~\cite{Wissel2007} expiries and strikes are not floating; they are fixed. The model will evolve
along the grid. Specifically, at time $t$ the future expiries in the model are those $T_j$ such that $T_j>t$. We assume that 
the largest strike spacing is $h:=\max_i(K^{i+1}-K^i)$.

\subsection{Discrete Local Volatility Backbone}

At time~$t$ we denote by $\dlv^j_t= (\dlv^{j,i}_t)_{i=1,\dots,N}$ the squared discrete local volatility at the
interior nodes $i=1,\dots,N$ for expiry $T_j > t$. We let $\lambda=\log\dlv$
throughout. Define
\begin{equation}\label{eq:weights}
  \omega^k:=\frac{(K^k)^{2}}{K^{k+1}-K^{k-1}},\qquad
  \iota^{k,-}:=\frac{\omega^k}{K^k-K^{k-1}},\qquad
  \iota^{k,+}:=\frac{\omega^k}{K^{k+1}-K^k},
  \qquad k=1,\dots,N,
\end{equation}
and then the tridiagonal matrix $A_j(\dlv)$ with rows
\begin{equation}\label{eq:operator}
  \bigl(A_j(\dlv)q\bigr)_i=
  \begin{cases}
    q^0-\Delta_j\dlv^1\iota^{1,-}q^1, & i=0,\\[2pt]
    -\Delta_j\dlv^{i-1}\iota^{i-1,+}q^{i-1}
     +\bigl(1+\Delta_j\dlv^i(\iota^{i,-}+\iota^{i,+})\bigr)q^i
     -\Delta_j\dlv^{i+1}\iota^{i+1,-}q^{i+1}, & 1\le i\le N,\\[2pt]
    -\Delta_j\dlv^N\iota^{N,+}q^N+q^{N+1}, & i=N+1 .
  \end{cases}
\end{equation}
(here $q^i$ denotes the $i$-th component of a generic $q\in\R^{N+2}$).
Up to notation this is the implicit forward operator of \cite{BuehlerRyskin2016}.

Let now $X_t$ be the current spot price at time~$t$ for the backbone process $X$ and let $q_t^0\in\R^{N+2}$
the density
$$
  q^{0,i}_t = \frac{  ( X_t - K^{i-1} )^+ - ( X_t - K^{i} )^+ }{ K^{i} - K^{i-1} } - \frac{ ( X_t - K^i )^+ - ( X_t - K^{i+1} )^+   }{ K^{i+1} - K^{i} } 
$$


Then, the family $q_t=(q_t^j)_{j=0,\dots,M}$ starting in $q_t^0$ is then given iteratively as
\begin{equation}
A_j(\dlv^j_t)\,q^j_t=q^{j-1}_t
\end{equation}
Note that in \eqref{eq:operator} the entry multiplying $q^{i\pm1}_t$ carries
$\dlv^{i\pm1}_t$ and not $\dlv^i_t$. The following was shown in~\cite{BuehlerRyskin2016}:

\begin{theorem}
The family is then $q=(q^{j,i}_t)$ with $q^j_t\in\R^{N+2}$ us a discrete risk-neutral density on the grid $(T,K)$
in the sense that
\begin{enumerate}
  \item $q^{j,i}_t \geq 0$
  \item $\sum_{i=0}^{N+1} q^{j,i}_t = 1$
  \item $\sum_{i=0}^{N+1} K^i q^{j,i}_t = X_0$.
\end{enumerate}
This means there is a discrete time-and-space process $X_t=(X_t^j)_{j=0,\dots,M}$ which takes values only in $(T,K)$. Note that
the model-time index here is~$j$.
Call prices for $X_t$ for arbitrary $(\tau,k)\in\R_{\geq 0}^2$ are given by the linear functional
\begin{equation}\label{eq:Xcall}
  C_t(k,\tau)=\mathbb{E}\left[ \left( X_t^j - k \right)^+ \right] = \sum_{i=0}^{N+1} ( K^i - k )^+ q^{j,i}_t \ .
\end{equation}
with $j=\max\{j: T_j\leq \tau\}$.
\end{theorem}








Because \eqref{eq:operator} is affine in $\dlv$ we may write
\begin{equation}\label{eq:affine}
  A_j(\dlv)\,q\;=\;q+G_j(q)\,\dlv,\qquad G_j(q)\in\R^{(N+2)\times N},
\end{equation}
so that $G_j(q)$ collects the terms carrying $\dlv$. The map $G_j$ is linear in
$q$ and $G_j(q)\dlv$ is linear in each argument separately; we use both
readings. The following elementary observation is used constantly.

\begin{lemma}[Factorisation]\label{lem:factor}
Column $k$ of $G_j(q)$ involves $q$ only through $q^k$. Consequently
\begin{equation}\label{eq:factor}
  G_j(q)\;=\;\widehat G_j\,\mathrm{diag}(q^1,\dots,q^N),
\end{equation}
where $\widehat G_j\in\R^{(N+2)\times N}$ has $k$-th column
$\Delta_j\bigl(-\iota^{k,-} e_{k-1}+(\iota^{k,-}+\iota^{k,+})e_k-\iota^{k,+} e_{k+1}\bigr)$.
Moreover $\widehat G_j$ has rank $N$, hence is injective.
\end{lemma}

\begin{proof}
Immediate from \eqref{eq:operator}, whose terms carrying $\dlv^k$ are all
proportional to $q^k$. For injectivity, suppose $\widehat G_jc=0$ and let $k$ be
the smallest index with $c_k\neq0$. Row $k-1$ receives a contribution from column
$k'$ only for $k'\in\{k-2,k-1,k\}$, and the first two carry $c_{k'}=0$. So row
$k-1$ reads $-\Delta_j\iota^{k,-}c_k=0$, forcing $c_k=0$, a contradiction.
\end{proof}

Note that $\widehat G_j$ depends only on the grid and on $\Delta_j$, not on the
state; the state enters \eqref{eq:factor} only through the diagonal factor.

\begin{proposition}[The code-book]\label{prop:codebook}
Fix the front density $q^{1}$, a probability vector on the grid with mean $S_0$, and let $j=2,\dots,M$.
\begin{enumerate}[label=(\roman*),leftmargin=2em]
\item Given densities $q^{j-1}$ and $q^{j}$ with $q^j>0$, the discrete local
variance is uniquely determined by the linear system
\begin{equation}\label{eq:bijection}
  G_j(q^j)\,\dlv^j\;=\;q^{j-1}-q^{j}.
\end{equation}
\item Conversely, given $\dlv^j>0$, the vector $q^j:=A_j(\dlv^j)^{-1}q^{j-1}$
is again a probability vector with mean $S_0$, and the gridded prices
$(\pi^i)^{\!\top}q^j$ admit no static arbitrage.
\item A gridded surface is free of static arbitrage, butterflies and
calendar spreads together, if and only if it arises as in \textup{(ii)} from
some $\dlv\ge0$, and it lies in the \emph{interior} of the arbitrage-free set
if and only if $\dlv>0$.
\end{enumerate}
\end{proposition}

\begin{proof}
(i) is \eqref{eq:affine}, with uniqueness from Lemma~\ref{lem:factor}. For (ii),
$A_j(\dlv^j)^{-1}$ is non-negative with unit column sums and inherits the
martingale property \cite{BuehlerRyskin2016}, attributed there
to Andreasen and Huge \cite{AndreasenHuge2011}; non-negativity excludes butterfly
arbitrage, and $q^j$ is then a mean-preserving spread of $q^{j-1}$, so Jensen's
inequality excludes calendar arbitrage. (iii) is \cite{BuehlerRyskin2016}, after
\cite{Buehler2006}: absence of arbitrage
is equivalent to non-negative Gamma and Backward-Theta at every node, and $\dlv$
is their ratio. Both conditions are
nodewise, and $\dlv^{j,i}=0$ exactly when Backward-Theta vanishes.
\end{proof}

What everything else rests on is that the admissible set is a product of
half-lines, one per node, so it has no faces in the sense that a polytope does,
and $\lambda=\log\dlv$ maps its interior diffeomorphically onto $\R^{N(M-1)}$.


\subsection{Two conservation laws}\label{sec:conservation}

The whole structure of the drift restriction follows from one observation about
the left null space of $G_j$. Both identities below are used by \cite{BuehlerRyskin2016}, where they hold by
construction for the backward operator. What is not there, because nothing in
that paper needs it, is that $\mathrm{span}\{\one,K\}$ is the \emph{entire} left
null space. That is a rank
statement, and it is what makes the drift system of Section~\ref{sec:drift}
square.

\begin{lemma}[Left null space]\label{lem:null}
For every $q$ and every $j$,
\[
  \one^{\!\top}G_j(q)=0
  \qquad\text{and}\qquad
  K^{\!\top}G_j(q)=0 ,
\]
where $\one=(1,\dots,1)$ and $K=(K^0,\dots,K^{N+1})$. Consequently $A_j(\dlv)$
conserves both total mass and the spot, for any $\dlv$. Combined with
Lemma~\ref{lem:factor}, the left null space of $\widehat G_j$ is exactly
$\mathrm{span}\{\one,K\}$.
\end{lemma}

\begin{proof}
Fix an interior index $k$ and collect the three entries of $G_j(q)$ in column
$k$: from \eqref{eq:operator} they are $\Delta_j(\iota^{k,-}+\iota^{k,+})q^k$ in row
$k$, $-\Delta_j\iota^{k,+}q^k$ in row $k+1$, and $-\Delta_j\iota^{k,-}q^k$ in row $k-1$.
Their sum is $0$, giving the first identity. Weighting by the node values,
\begin{align*}
  \Delta_jq^k\Bigl[K^k\bigl(\iota^{k,-}+\iota^{k,+}\bigr)
   -K^{k+1}\iota^{k,+}-K^{k-1}\iota^{k,-}\Bigr]
  &=\Delta_jq^k\Bigl[\iota^{k,-}(K^k-K^{k-1})-\iota^{k,+}(K^{k+1}-K^k)\Bigr]\\
  &=\Delta_jq^k\bigl[\omega^k-\omega^k\bigr]=0
\end{align*}
by \eqref{eq:weights}, giving the second. The conservation statements follow
from \eqref{eq:affine}: $\one^{\!\top}A_j(\dlv)q=\one^{\!\top}q$ and
$K^{\!\top}A_j(\dlv)q=K^{\!\top}q$. Since $\widehat G_j$ has rank $N$ and lives
in $\R^{(N+2)\times N}$, its left null space is two-dimensional, hence exactly
$\mathrm{span}\{\one,K\}$.
\end{proof}

\begin{remark}[Counting]\label{cor:counting}
For $q>0$, $G_j(q)$ has $N+2$ rows and $N$ columns and rank $N$. Hence the $N+2$ martingale
conditions per expiry contain exactly two redundancies, conservation of mass
and of the spot, and the remaining $N$ conditions match the $N$ unknowns
of $\dlv^j$. The system is square on the quotient.
\end{remark}

\section{The drift restriction}\label{sec:drift}

\begin{assumption}[Martingale property]\label{ass:mart}
Under $\Q$ every gridded instrument price is a martingale in $t$. Since the
gridded payoffs $\pi^i$ span $\R^{N+2}$, this is equivalent to each density
vector being a martingale,
\begin{equation}\label{eq:qdyn}
  \dd q^j_t=\psi^j_t\,\dd W_t,\qquad \psi^j_t\in\R^{(N+2)\times d}.
\end{equation}
\end{assumption}

This is the only assumption behind the drift restriction, and it carries no
regularity. Each
$q^j$ is a fixed linear combination of traded butterflies, so \eqref{eq:qdyn} is
the martingale representation theorem applied to traded prices, and nothing has
to commute with a stochastic integral. This is the substantive advantage of a
discrete code-book over its continuum analogue, where the corresponding step
requires differentiating a stochastic field in strike and maturity and must be
licensed by standing regularity conditions. Theorem~\ref{thm:pointwise} below
needs only $q^j>0$.

\begin{theorem}[Drift restriction]\label{thm:drift}
Let $\dd\dlv^j_t=\alpha^j_t\dd t+\beta^j_t\dd W_t$. Under
Assumption~\ref{ass:mart}, for each $j=1,\dots,M$,
\begin{align}
  \text{(volatility)}\qquad
  &A_j(\dlv^j)\,\psi^j\;=\;\psi^{j-1}-G_j(q^j)\,\beta^j,
  \label{eq:volsolve}\\[4pt]
  \text{(drift)}\qquad
  &G_j(q^j)\,\alpha^j\;=\;-\sum_{f=1}^{d}G_j\bigl(\psi^{j,f}\bigr)\beta^{j,f},
  \label{eq:driftsolve}
\end{align}
marched forward in $j$ from $\psi^0$, the volatility of the initial density,
which is determined by the spot dynamics.
\end{theorem}

\begin{proof}
Apply It\^o to $A_j(\dlv^j_t)q^j_t=q^{j-1}_t$, using \eqref{eq:affine} to write
the left side as $q^j+G_j(q^j)\dlv^j$ and noting that $G_j$ is bilinear. The
$\dd W$ terms give $\psi^j+G_j(\psi^j)\dlv^j+G_j(q^j)\beta^j=\psi^{j-1}$, which
is \eqref{eq:volsolve} by \eqref{eq:affine}. The $\dd t$ terms give
$G_j(q^j)\alpha^j+\sum_f G_j(\psi^{j,f})\beta^{j,f}=0$, since $q^{j-1}$ and
$q^{j}$ are driftless by Assumption~\ref{ass:mart} and the only quadratic
covariation is between $G_j(\cdot)$ and $\dlv^j$, both driven by $W$.
\end{proof}

Equation \eqref{eq:volsolve} is not ours. Compared with the pricing recursion
$A_j(\dlv^j)q^j=q^{j-1}$ it is \emph{the same operator with one added source},
the structure of Carmona and Nadtochiy \cite{CarmonaNadtochiy2008}. The drift
restriction below is theirs in the continuum, and Wissel's \cite{Wissel2007} on
this lattice; we claim no priority for either.

What the fixed lattice changes is threefold. Their restriction carries a term in
the derivative of the code-book along time to maturity; a lattice of fixed expiry
dates has no such derivative, and here that term is absent. The relation is algebraic rather than
differential, so neither the high-order regularity their statement requires nor
the singular behaviour as time to maturity vanishes that it is built to
accommodate can arise here, the front slab having positive width. And the
discrete system is over-determined, $N+2$ conditions on $N$ unknowns, which is
what the counting of Remark~\ref{cor:counting} squares; that counting is what
turns a differential relation into an identity, and it has no continuum
counterpart. Appendix~\ref{app:priority} carries both identifications in the
original notation.

\begin{theorem}[The drift is a pointwise covariance]\label{thm:pointwise}
Define the relative density volatility
\begin{equation}\label{eq:eta}
  \eta^{j,i}\;:=\;\frac{\psi^{j,i}}{q^{j,i}},
  \qquad i=0,\dots,N+1,
\end{equation}
that is, the relative volatility of $q^{j,i}$. Only the interior components enter
\eqref{eq:pointwise}, since $\dlv$ lives on the interior nodes; the boundary
components are needed in Section~\ref{sec:exist}. Then \eqref{eq:driftsolve} is solvable
for every $\beta$ and every state, its solution is unique, and it is given
componentwise by
\begin{equation}\label{eq:pointwise}
  \alpha^{j,i}\;=\;-\ip{\eta^{j,i}}{\beta^{j,i}}
  \qquad\Longleftrightarrow\qquad
  \dd\lambda^{j,i}=\Bigl[-\ip{\eta^{j,i}}{b^{j,i}}-\tfrac12\abs{b^{j,i}}^{2}\Bigr]\dd t
   +b^{j,i}\,\dd W_t,
\end{equation}
where $\lambda=\log\dlv$ and $b=\beta/\dlv$.
\end{theorem}

\begin{proof}
By Lemma~\ref{lem:factor} both sides of \eqref{eq:driftsolve} factor through the
same fixed matrix:
\[
  \widehat G_j\,\mathrm{diag}(q^j)\,\alpha^j
  \;=\;-\sum_f\widehat G_j\,\mathrm{diag}(\psi^{j,f})\,\beta^{j,f} .
\]
Since $\widehat G_j$ is injective it cancels, leaving
$\mathrm{diag}(q^j)\alpha^j=-\sum_f\mathrm{diag}(\psi^{j,f})\beta^{j,f}$, which
is \eqref{eq:pointwise} after dividing by $q^{j,i}>0$. Existence and uniqueness
follow from the same cancellation. The logarithmic form is It\^o applied to
$\lambda=\log\dlv$, using $\alpha/\dlv=-\ip{\eta}{b}$.
\end{proof}

Two features of \eqref{eq:pointwise} are worth drawing out. The first is its
shape. It is the exact analogue of the HJM restriction $\alpha=\sigma\int\sigma$:
the drift of the log discrete local variance is minus the covariance of its own
volatility with the volatility of the log density, plus an It\^o convexity term.
Where HJM's derived object is an integral of the volatility across maturities,
ours is $\eta$, obtained from $b$ by the single forward recursion
\eqref{eq:volsolve}. The second is the cost. Once $\eta$ is known no matrix is inverted: no Jacobian
of the pricing map is ever formed and no rank condition on it appears. And since
$\psi$ is linear in $\beta$ and $\alpha$ is bilinear in $(\beta,\psi)$, the drift
is an exact quadratic form in $\beta$, as in HJM; Theorem~\ref{thm:wing} is what
rules out the corresponding finite-time explosion.

The driftless point is the classical model: if $\beta\equiv0$ and $\psi^0=0$ then
$\psi\equiv0$ and $\alpha\equiv0$, so the discrete local variance surface is
frozen, and with $\psi^0\ne0$ the densities still move, driven by the spot alone,
with $\alpha\equiv0$ persisting.

\begin{proposition}[Spot consistency]\label{prop:spot}
The spot is $S_t=K^{\!\top}q^{j}_t$, the same for every $j$ by
Lemma~\ref{lem:null}. Since $K^{\!\top}A_j=K^{\!\top}$ and
$K^{\!\top}\widehat G_j=0$, equation \eqref{eq:volsolve} gives
\[
  \one^{\!\top}\psi^{j}=\one^{\!\top}\psi^{0}=0
  \qquad\text{and}\qquad
  K^{\!\top}\psi^{j}=K^{\!\top}\psi^{0}=\sigma^{S}_tS_t
  \quad\text{for every }j.
\]
So the spot's volatility is the single scalar $K^{\!\top}\psi^{0}$, it is
automatically the same at every expiry, and no further condition links the spot
to the surface. Taking $\psi^{0}=0$ freezes the spot and recovers
$K^{\!\top}\psi^{j}=0$.
\end{proposition}

Proposition~\ref{prop:spot} is the discrete form of Carmona and Nadtochiy's spot
consistency condition: the spot and its volatility are both \emph{read off} the
surface rather than modelled, $S=K^{\!\top}q^j$ and $\sigma^SS=K^{\!\top}\psi^j$
for any $j$.

The drift restriction is used below through the following chain, which is what a
model actually evaluates:
\begin{equation}\label{eq:pipeline}
  b\;\xrightarrow{\ \eqref{eq:volsolve}\ }\;\psi^1,\dots,\psi^M
  \;\xrightarrow{\ \eta=\psi/q\ }\;\alpha ,
\end{equation}
a forward march of $d$ tridiagonal solves per expiry, $O(NMd)$ in total. The
volatility structure is the only input; the drift is computed, never fitted.

\section{The spot and the front expiry}\label{sec:front}

The spot is not a state variable. By Proposition~\ref{prop:spot} it is the read-off
$S_t=K^{\!\top}q^j_t$, the same scalar at every expiry. What has to be specified is
the front density $q^{1}$, the initial condition of the recursion of
Proposition~\ref{prop:codebook}; everything at $j\ge2$ follows from it and from the
volatility structure.

\subsection{The front chart}\label{sec:frontslice}

The front density must be strictly positive and sum to one. Both are automatic in the
chart
\begin{equation}\label{eq:softmax}
  q^{1,i}=\frac{e^{z^i}}{\sum_ke^{z^k}},\qquad
  \dd z=m\,\dd t+\sum_{f=1}^{d}\nu^f\,\dd W^f ,
\end{equation}
with $z\in\R^{N+2}$ unconstrained. The chart is invariant under $z\mapsto z+c\one$, so a
shift of $\nu$ along $\one$ leaves $q^1$ unchanged and only the $q^1$-centred part of
$\nu$ is identified; write $\mathrm{dev}^f:=\nu^f-(q^1)^{\!\top}\nu^f\,\one$ for it. The input
$\nu$ is free, and plays the role $b$ plays at $j\ge2$. It\^o's formula for a function
of a vector process shows that each $q^{1,i}$ is a martingale exactly when
\begin{equation}\label{eq:frontdrift}
  m^i=-\tfrac12\sum_f\bigl(\mathrm{dev}^{i,f}\bigr)^2
\end{equation}
up to a multiple of $\one$. Furthermore we obtain
\begin{equation}\label{eq:psi1}
  \psi^{1,\cdot,f}=q^1\circ\mathrm{dev}^f,
  \qquad
  \eta^{1,\cdot,f}=\mathrm{dev}^f,
  \qquad
  \one^{\!\top}\psi^{1,\cdot,f}=0 ,
\end{equation}
with $\circ$ the entrywise product.
The same computation appears in \cite{FilipovicHughstonMacrina2012} and as the
stochastic replicator equation \cite{MertikopoulosMoustakas2010}.

\begin{remark}[Resolution at the front expiry]\label{rem:frontres}
Among grid measures with mean $S$ the least dispersed is the pair of nodes straddling
it, so
\begin{equation}\label{eq:varfloor}
  \Var_{q^1}(K)\;\ge\;(S-K^i)(K^{i+1}-S),
  \qquad K^i\le S<K^{i+1},
\end{equation}
which is at most $\tfrac14h^{2}$. The front implied volatility therefore has the floor
\[
  \frac{\sqrt{(S-K^i)(K^{i+1}-S)}}{S\sqrt{T_1-t}}\;\le\;\frac{h}{2S\sqrt{T_1-t}},
\]
largest at the midpoint of a cell and zero when the spot sits on a node, which is
exactly what Proposition~\ref{prop:collapse} forces at $T_1$.
\end{remark}

\subsection{Terminal settlement: pinning the spot to a strike}\label{sec:nonpinning}

The chart constrains $q^1$ at no time $t<T_1$ beyond positivity and mass, so it imposes
no terminal condition. Wissel's front parametrisation fixes the terminal state whatever
the coefficients \cite{Wissel2007}; here the terminal state is a modelling choice. We take
the front volatility proportional to the distance from the spot, in one factor,
\begin{equation}\label{eq:nufront}
  \eta^{1,i}_{t}=\frac{c_t\,(K^i-S_t)}{h},\qquad c_t=a_1\tau^{-p}\ \text{near }T_1,
  \quad\tau=T_1-t,\ a_1>0 .
\end{equation}
Only the behaviour near $T_1$ matters: on $[0,T_1-\varepsilon]$ the front
volatility may be any bounded, locally Lipschitz choice, and \eqref{eq:nufront} need only
hold on $[T_1-\varepsilon,T_1)$. It is already centred, since $\sum_iq^{1,i}(K^i-S_t)=0$. Its coefficients are polynomial in
$q^1$ and bounded on $[0,T]$ for every $T<T_1$, so there is a unique strong solution
on $[0,T_1)$,
\begin{equation}\label{eq:frontexplicit}
  q^{1,i}_{t}=q^{1,i}_{0}\exp\Bigl(\int_0^t\eta^{1,i}\,\dd W-\tfrac12\int_0^t(\eta^{1,i})^2\dd s\Bigr),
  \qquad
  \lVert\eta^{1}_t\rVert_\infty\le c_t\,w/h,\quad w:=K^{N+1}-K^0 .
\end{equation}

\begin{proposition}[The front collapses onto a node, and settles exactly]\label{prop:collapse}
If $p\ge\tfrac12$ then $q^{1}_{T_1}:=\lim_{t\uparrow T_1}q^1_t=e_k$ for a random index $k$,
almost surely, so $S_{T_1}=K^k$ and
\begin{equation}\label{eq:settle}
  \sum_i q^{1,i}_{T_1}(K^i-K^n)^{+}=(S_{T_1}-K^n)^{+}\qquad\text{for every }n .
\end{equation}
If $p<1$ then Assumption~\ref{ass:F} holds on $[0,T_1]$. The window $p\in[\tfrac12,1)$ gives both.
\end{proposition}

\begin{proof}
Each $q^{1,i}$ is a bounded martingale, so $q^1_t\to q^1_{T_1}$ almost surely, and the
variance $V_t:=\sum_iq^{1,i}_{t}(K^i-S_t)^2$ converges to that of $q^1_{T_1}$. By
\eqref{eq:psi1} and \eqref{eq:nufront} the volatility of $S=K^{\!\top}q^1$ is
$\sum_iK^iq^{1,i}\eta^{1,i}=c_tV_t/h$ on $[T_1-\varepsilon,T_1)$. Since $S$ is a martingale with values in
$[K^0,K^{N+1}]$,
\[
  \E\int_{T_1-\varepsilon}^{T_1}\frac{c_t^2V_t^2}{h^2}\,\dd t
  =\E\bigl[(S_{T_1}-S_{T_1-\varepsilon})^2\bigr]\le w^2 ,
\]
so the integral is finite almost surely. On the event $\{V_{T_1}>0\}$ we would have
$V_t\ge V_{T_1}/2$ near $T_1$, and the integral would diverge, because
$\int^{T_1}c_t^2\,\dd t=a_1^2\int^{T_1}\tau^{-2p}\,\dd t=\infty$ for $2p\ge1$. Hence
$V_{T_1}=0$ almost surely: $q^1_{T_1}$ is a probability vector on distinct points with
zero variance, that is $e_k$ with $K^k=S_{T_1}$, and \eqref{eq:settle} reads
$(K^k-K^n)^{+}=(S_{T_1}-K^n)^{+}$. Finally $\eta^1$ is bounded on $[0,T_1-\varepsilon]$ and
$\int_{T_1-\varepsilon}^{T_1}\lVert\eta^1_t\rVert_\infty\dd t\le
(w/h)\int_{T_1-\varepsilon}^{T_1}c_t\,\dd t<\infty$ when $p<1$, by
\eqref{eq:frontexplicit}.
\end{proof}

The terminal spot is therefore a grid point, as in \cite{Wissel2007}; the difference is
that here it is chosen, by \eqref{eq:nufront}, and there it is forced by the
parametrisation.

\begin{remark}[What the lattice could carry, and does not]\label{rem:bary}
At expiry each butterfly price $q^{j,i}$ is worth its payoff: for $S_{T_j}\in[K^k,K^{k+1}]$ the
second difference defining $q^{j,i}$ telescopes to $\varpi(S_{T_j})$, the barycentric
weights of the spot on the two nodes straddling it. Since $u\mapsto(u-K^n)^{+}$ is affine
between nodes,
\begin{equation}\label{eq:barypayoff}
  \sum_i \varpi^i(s)\,(K^i-K^n)^{+}\;=\;(s-K^n)^{+}
  \qquad\text{for every }s\in[K^0,K^{N+1}]\text{ and every model strike }K^n ,
\end{equation}
so $\varpi(s)$ is an admissible terminal state that settles every gridded contract exactly,
also for $s$ strictly inside a cell. The chart does not force the pinning; the choice
\eqref{eq:nufront} does, because its terminal variance must vanish. A front volatility
that drives $q^1$ to $\varpi(s)$ with $s$ inside a cell is left as future work
(Section~\ref{sec:open}).
\end{remark}

\section{The wing constant, and existence}\label{sec:exist}

The log discrete local variance $\lambda=(\lambda^2,\dots,\lambda^M)$,
$\lambda^j=\log\dlv^j$, ranges over all of $\R^{N(M-1)}$, so positivity is
automatic; the front density has its own chart \eqref{eq:softmax}. By Theorem~\ref{thm:pointwise} the model
is the stochastic differential equation
\begin{equation}\label{eq:abstract}
  \dd\lambda_t=\Bigl[-\ip{\eta_t}{b_t}-\tfrac12\abs{b_t}^{2}\Bigr]\dd t
   +b_t\,\dd W_t ,
\end{equation}
whose only non-elementary coefficient is $\eta=\psi/q$ from \eqref{eq:volsolve}
and \eqref{eq:eta}. Its denominator is small in the wings, and the question is
whether $\eta$ can blow up there.

The argument has the structure of Wissel's proof of his Theorem~3.5
\cite{Wissel2007}: log coordinates, induction over expiries, and local existence
followed by a bound that rules out explosion. Its one non-standard ingredient is an a priori bound, independent of the state, on the
relative volatility of the calendar spread at the second expiry. Wissel could not
obtain this bound in general and imposes his condition (3.24) instead (his
Remark~1 following the proof, p.~18); in our coordinates the quantity is $\eta+b$
(Appendix~\ref{app:wissel}). We obtain it from an exact identity for the
volatility recursion and a known inequality for inverses of M-matrices.

\subsection{The volatility recursion}\label{sec:walk}\label{sec:master}

Write $P:=A_j(\dlv^j)^{-1}$, so that $q^j=Pq^{j-1}$.

\begin{corollary}[Positivity, unit column sums, absorbing ends]\label{cor:pos}
$P\ge0$ entrywise, $\one^{\!\top}P=\one^{\!\top}$, so $P_{im}\in[0,1]$, and the outer
columns of $P$ are $e_0$ and $e_{N+1}$.
\end{corollary}

\begin{proof}
By \eqref{eq:operator}, $A_j$ has non-positive off-diagonal entries, and in each
column the diagonal entry exceeds the sum of the moduli of the off-diagonal ones by
exactly one; so $A_j$ is a non-singular M-matrix and $P\ge0$
\cite{BermanPlemmons1994}. The column sums follow from $\one^{\!\top}A_j=\one^{\!\top}$
(Lemma~\ref{lem:null}), and the outer columns of $A_j$ are $e_0$ and $e_{N+1}$. The
first two statements are also Theorem~3.3 of \cite{BuehlerRyskin2016}, attributed
there to \cite{AndreasenHuge2011}.
\end{proof}

That is all that is needed about $P$ itself. What we want next is a recursion for
$\eta$ alone: one step per expiry, each adding a controlled amount, so that the
bound telescopes along the ladder. Dividing \eqref{eq:volsolve} by $q^j$ gives
\begin{equation}\label{eq:divided}
  \eta^{j,i}\;=\;\frac{(P\psi^{j-1})_i}{q^{j,i}}
   \;-\;\frac{\bigl(P\,G_j(q^j)\beta^j\bigr)_i}{q^{j,i}} ,
\end{equation}
the first term carrying $\eta^{j-1}$ and the second carrying $b^j$.

The first is an average. Writing $\psi^{j-1}=q^{j-1}\!\circ\eta^{j-1}$, its weights
$P_{ik}q^{j-1,k}/q^{j,i}$ are non-negative because $P\ge0$ and sum to one because
$q^j=Pq^{j-1}$, so it contributes at most $\lVert\eta^{j-1}\rVert_\infty$; they are
the matrix $\mathcal A$ of \eqref{eq:mastercomp}.

The source is the difficulty. Since $G_j$ is a second difference, row $i$ carries
the neighbouring densities $q^{j,i\pm1}$ as well as $q^{j,i}$, so dividing by $q^{j,i}$
leaves the ratios
\begin{equation}\label{eq:nbr}
  \frac{q^{j,i-1}}{q^{j,i}}
  \qquad\text{and}\qquad
  \frac{q^{j,i+1}}{q^{j,i}} ,
\end{equation}
which are unbounded in the wings, where the density decays fastest; nothing in
\eqref{eq:volsolve} bounds them.

They are, however, an artefact of writing the row out. Since
$\one^{\!\top}G_j(q)=0$ by Lemma~\ref{lem:null}, row $i$ of $G_j(q)\dlv^j$ is the
net outflow of node $i$, and by \eqref{eq:affine} that net is
\begin{equation}\label{eq:netknown}
  G_j(q^j)\,\dlv^j\;=\;A_j(\dlv^j)q^j-q^j\;=\;q^{j-1}-q^{j},
\end{equation}
which carries no ratios at all: the individual fluxes are wing quantities, their
net is not. For $b^j\equiv c$ this settles the matter, since $\beta^j=c\,\dlv^j$
and $Pq^{j-1}=q^j$ give
\[
  \frac{\bigl(P\,G_j(q^j)\beta^j\bigr)_i}{q^{j,i}}
  =c\,\frac{\bigl(q^{j}-Pq^{j}\bigr)_i}{q^{j,i}}
  =c\Bigl(1-\frac{(Pq^{j})_i}{q^{j,i}}\Bigr),
\]
the two ratios of \eqref{eq:nbr} having collapsed into the single propagation
ratio $(Pq^j)_i/q^{j,i}$. A general $b$ only modulates the density. Let
$E\in\R^{N\times(N+2)}$ restrict to the interior nodes, so $EE^{\!\top}=I_N$, and
note that $A_j=I+\Gh_j\,\mathrm{diag}(\dlv^j)E$. Writing
$E^{\!\top}(q\circ b)$ for the vector equal to $q^ib^i$ at the interior nodes and
zero at the outer ones, \eqref{eq:factor} and the commuting of diagonal matrices
give
\begin{equation}\label{eq:modulated}
  G_j(q)\beta^j
  =\Gh_j\,\mathrm{diag}(\dlv^j)E\,E^{\!\top}(q\circ b^j)
  =\bigl(A_j(\dlv^j)-I\bigr)E^{\!\top}(q\circ b^j),
\end{equation}
again the operator's own increment, applied to $E^{\!\top}(q\circ b^j)$ in place of
$q$. The proposition is what this amounts to.

\begin{proposition}\label{prop:master}
With $b^j=\beta^j/\dlv^j$ and $\circ$ the entrywise product,
\begin{equation}\label{eq:master}
  \psi^{j}\;=\;P\,\psi^{j-1}\;-\;(I-P)\,E^{\!\top}\bigl(q^{j}\!\circ b^{j}\bigr),
\end{equation}
equivalently, with the entrywise non-negative
\[
  \mathcal A:=\mathrm{diag}(q^{j})^{-1}P\,\mathrm{diag}(q^{j-1}),
  \qquad
  \mathcal B:=\mathrm{diag}(q^{j})^{-1}P\,E^{\!\top}\mathrm{diag}(Eq^{j}),
\]
\begin{equation}\label{eq:mastercomp}
  \eta^{j}\;=\;\mathcal A\,\eta^{j-1}\;+\;\bigl(\mathcal B-E^{\!\top}\bigr)b^{j} ,
\end{equation}
where $\mathcal A\one=\one$, so the transport is an averaging operator, and the
$i$-th row of $\mathcal B$ sums to
$(Pq^{j})_i/q^{j,i}-\one_{\{i\in\{0,N+1\}\}}$.
\end{proposition}

\begin{proof}
By \eqref{eq:volsolve}, $\psi^j=P\psi^{j-1}-P\,G_j(q^j)\beta^j$. Applying $P$ to
\eqref{eq:modulated} at $q=q^j$ and using $PA_j=I$ gives
$P\,G_j(q^j)\beta^j=(I-P)E^{\!\top}(q^j\!\circ b^j)$, which is \eqref{eq:master}.
Dividing it by $q^{j,i}$ gives \eqref{eq:mastercomp}: the rows of $\mathcal A$ sum to
$(Pq^{j-1})_i/q^{j,i}=1$, and for $\mathcal B$ the terms $k\in\{0,N+1\}$ missing from
$(Pq^j)_i$ vanish for interior $i$ and equal $q^{j,i}$ at $k=i$ for outer $i$, by
Corollary~\ref{cor:pos}.
\end{proof}

So $\eta^j$ is an average of $\eta^{j-1}$ plus a source in $b^j$, and the march
across expiries is additive.

\subsection{The wing constant}\label{sec:bound}

That single surviving ratio, $(Pq^j)_i/q^{j,i}$, is what has to be bounded. Since
$q^j=Pq^{j-1}$ it is
\begin{equation}\label{eq:P2ratio}
  \frac{(Pq^{j})_i}{q^{j,i}}=\frac{(P^{2}q^{j-1})_i}{(Pq^{j-1})_i} ,
\end{equation}
so what is needed is how much larger than $P$ the square $P^{2}$ can be entrywise;
and since both sides are linear and non-negative in $q^{j-1}$, the case
$q^{j-1}=e_m$ shows that an entrywise bound is not merely sufficient but
necessary. Write
\begin{equation}\label{eq:rows}
  R^{j,i}:=\one_{\{1\le i\le N\}}+\frac{(Pq^j)_i}{q^{j,i}}-\one_{\{i\in\{0,N+1\}\}},
  \qquad
  \Theta_j:=\max_i R^{j,i},
  \qquad
  \Xi:=\sum_{r=1}^{N}P_{rr}\le N ,
\end{equation}
and $H_j:=\max_i\abs{\eta^{j,i}}$. Here $R^{j,i}$ is the row bound of
\eqref{eq:mastercomp} at node $i$, and $\Theta_j$, the \emph{wing constant}, is its
size across the slab. The theorem shows $\Theta_j$ is bounded by the grid alone,
uniformly in the state; the bound rests on a single inequality for inverses of
M-matrices.

\begin{lemma}[Squared propagator]\label{lem:pathproduct}\label{cor:p2p}
For all nodes $i,r,m$,
\begin{equation}\label{eq:summand}
  P_{ir}\,P_{rm}\;\le\;P_{im}\,P_{rr},
\end{equation}
and consequently, for all $i,m$,
\begin{equation}\label{eq:p2p}
  \bigl(P^{2}\bigr)_{im}\le
  \begin{cases}
    \Xi\,P_{im}, & 1\le i\le N,\\
    (1+\Xi)\,P_{im}, & i\in\{0,N+1\}.
  \end{cases}
\end{equation}
\end{lemma}

\begin{proof}
$P$ is the inverse of an M-matrix by Corollary~\ref{cor:pos}, and every inverse
M-matrix is a path product matrix, $a_{ij}a_{jk}\le a_{ik}a_{jj}$ for all
$i,j,k$ \cite[Theorem~4.5]{JohnsonSmith2011}; \eqref{eq:summand} is the case
$(i,j,k)=(i,r,m)$. Summing it over $r$ gives \eqref{eq:p2p}: an outer $r$
contributes only if $r=i$, since column $r$ of $P$ is $e_r$, and then contributes
$P_{im}$.
\end{proof}

\begin{theorem}\label{thm:wing}
For every $\dlv\in(0,\infty)^N$,
\begin{equation}\label{eq:step}
  H_j\;\le\;H_{j-1}+\Theta_j\lVert b^j\rVert_\infty ,
\end{equation}
and
\begin{equation}\label{eq:thetabound}
  \Theta_j\le1+\Xi\le N+1,
  \qquad
  \Theta^{*}:=\sup_{\dlv}\sum_{j=2}^{M}\Theta_j\le(M-1)(N+1).
\end{equation}
\end{theorem}

\begin{proof}
Since $q^j=Pq^{j-1}\ge0$, Lemma~\ref{lem:pathproduct} gives
$(Pq^j)_i=\sum_m(P^2)_{im}q^{j-1,m}\le\Xi\,q^{j,i}$ for interior $i$ and
$\le(1+\Xi)q^{j,i}$ for outer $i$, so $R^{j,i}\le1+\Xi$ in both cases. Finally,
$\mathcal A\ge0$ with $\mathcal A\one=\one$ gives
$\lVert\mathcal A\eta^{j-1}\rVert_\infty\le H_{j-1}$, while row $i$ of
$(\mathcal B-E^{\!\top})b^{j}$ has modulus at most $\lVert b^j\rVert_\infty R^{j,i}$
by the row sums of Proposition~\ref{prop:master}.
\end{proof}

\subsection{Existence}\label{sec:existence}

Iterating \eqref{eq:step} from the front slice,
\begin{equation}\label{eq:apriori}
  \sup_{j,i}\abs{\eta^{j,i}_{t}}
  \;\le\;\lVert\eta^1_t\rVert_\infty+(M-1)(N+1)\,\lVert b\rVert_\infty ,
\end{equation}
and iterating \eqref{eq:master} shows that $\eta^j$ is affine in $\eta^1$, with
coefficients depending on $\lambda$ only through $P_2,\dots,P_j$ and $q^2,\dots,q^j$.
The only $t$ on the right of \eqref{eq:apriori} is inside $\eta^1_t$, which enters
linearly, so it has to be integrable in time rather than bounded.

\begin{assumption}[Integrable front volatility]\label{ass:F}
For a horizon $T\le T_1$,
\begin{equation}\label{eq:F}
  \int_0^{T}\lVert\eta^{1}_t\rVert_\infty\,\dd t\;<\;\infty\qquad\text{almost surely.}
\end{equation}
\end{assumption}

\begin{theorem}[Existence]\label{thm:exist}
Let $T<T_1$, and let the front volatility $\nu=\nu(t,z)$ in \eqref{eq:softmax} and
the volatility $b=b(t,z,\lambda)$ be bounded, $\abs{\nu}\le B_\nu$,
$\abs{b}\le B$, and locally Lipschitz in $(t,z,\lambda)$. Then the front chart
\eqref{eq:softmax}--\eqref{eq:frontdrift} together with \eqref{eq:abstract} has a
unique strong solution $(z,\lambda)$ on $[0,T]$, and
\begin{equation}\label{eq:ellbounds}
  0\;<\;\inf_{t\le T}\min_{j,i}\dlv^{j,i}_{t}
  \;\le\;\sup_{t\le T}\max_{j,i}\dlv^{j,i}_{t}\;<\;\infty
  \qquad\text{almost surely.}
\end{equation}
No condition is imposed on $\dlv$, on the wings, or on the number of factors.
\end{theorem}

\begin{proof}
The coefficients are locally Lipschitz: $q^1=\mathrm{softmax}(z)$ and
$\dlv=e^\lambda$ are positive, so each $A_j$ is a non-singular M-matrix and
everything except $\nu$ and $b$ is smooth in $(z,\lambda)$. Adding time as a state
variable to cover the time dependence, \cite[Theorem~V.38, p.~303]{Protter2005}
gives a unique solution up to an explosion time $\zeta$, with
$\limsup_{t\to\zeta}\lVert(z_t,\lambda_t)\rVert=\infty$ on $\{\zeta<\infty\}$.
The drift of $\lambda$ is bounded by \eqref{eq:apriori}, since
$\lVert\eta^1\rVert_\infty=\lVert\mathrm{dev}\rVert_\infty\le2B_\nu$, and the drift
$m$ of $z$ is bounded by $2B_\nu^2$; let $D$ bound both. On $[0,\zeta\wedge T)$
the drift part of $(z,\lambda)$ is therefore at most $DT$, and the noise part,
having integrand bounded by $B_\nu+B$, is a martingale bounded in $L^2$, hence
convergent as $t\to\zeta\wedge T$. So $(z,\lambda)$ stays bounded on
$[0,\zeta\wedge T)$, which rules out $\zeta\le T$. The solution exists on $[0,T]$,
$\lambda$ is bounded there, and \eqref{eq:ellbounds} follows from $\dlv=e^\lambda$.
\end{proof}

\begin{remark}[The front expiry is included]\label{rem:horizon}
At $T_1$ the front settles on a node (Proposition~\ref{prop:collapse}), so $z$
itself runs off and Theorem~\ref{thm:exist} does not apply on $[0,T_1]$ directly.
It applies on every $[0,T']$ with $T'<T_1$, and the $\lambda$-part of the proof
only needs the drift bound to be integrable in time, which is
Assumption~\ref{ass:F} with $T=T_1$. Under it $\lambda$ stays bounded on $[0,T_1)$
and has a limit at $T_1$, so \eqref{eq:ellbounds} holds with $T=T_1$. The settling
front has $\lVert\eta^1_t\rVert_\infty\le C(T_1-t)^{-p}$, integrable for $p<1$, so
the code-book reaches its first expiry with every gridded contract settling at its
payoff.
\end{remark}

\subsection{Beyond the first expiry}\label{sec:patching}

At $T_1$ the front expires and the next expiry becomes the front. The settling of
Proposition~\ref{prop:collapse} makes this switch clean: the spot sits on a node, and the
new front density is a column of the propagator. Indeed, if $q^l_{T_l}=e_k$ then the
code-book $A_{l+1}(\dlv^{l+1})\,q^{l+1}=q^l$ gives
\begin{equation}\label{eq:handover}
  q^{l+1}_{T_l}=P_{l+1}\,e_k ,\qquad P_{l+1}:=A_{l+1}\bigl(\dlv^{l+1}_{T_l}\bigr)^{-1},
\end{equation}
the spread of a unit mass at the settled spot by one step of the code-book. To restart the
front chart we need this density to be strictly positive.

\begin{lemma}[The new front]\label{lem:newfront}
Let $P=A_j(\dlv^j)^{-1}$ with $\dlv^j\in(0,\infty)^N$. If $k$ is an interior node, then
$(Pe_k)^i>0$ for every node $i$ and $K^{\!\top}Pe_k=K^k$. If $k\in\{0,N+1\}$, then
$Pe_k=e_k$.
\end{lemma}

\begin{proof}
Order the nodes with the interior first. The outer columns of $A_j$ are $e_0$ and
$e_{N+1}$ (Corollary~\ref{cor:pos}), so
\[
  A_j=\begin{pmatrix}A_{II}&0\\A_{BI}&I\end{pmatrix},\qquad
  P=\begin{pmatrix}A_{II}^{-1}&0\\-A_{BI}A_{II}^{-1}&I\end{pmatrix}.
\]
By \eqref{eq:operator}, $A_{II}$ is tridiagonal with strictly negative off-diagonal
entries, hence an irreducible non-singular M-matrix, so $A_{II}^{-1}>0$ entrywise
\cite{BermanPlemmons1994}. The block $-A_{BI}\ge0$ has a positive entry in the row of
$K^0$ (column $K^1$) and in the row of $K^{N+1}$ (column $K^N$), so the outer entries of an
interior column of $P$ are positive as well. The identity $K^{\!\top}P=K^{\!\top}$ follows
from $K^{\!\top}A_j=K^{\!\top}$ (Lemma~\ref{lem:null}), and the outer columns are
Corollary~\ref{cor:pos}.
\end{proof}

\begin{corollary}[Existence on the whole ladder]\label{cor:patch}
Let $b$ be bounded and locally Lipschitz. On each interval $[T_{l-1},T_l]$,
$l=1,\dots,M$, let the current front density $q^l$ follow the chart \eqref{eq:softmax}
with the settling volatility \eqref{eq:nufront}, where $\tau=T_l-t$ and
$p\in[\tfrac12,1)$. Then the model has a solution on $[0,T_M]$, unique given these
choices, such that
\begin{enumerate}[label=(\roman*),leftmargin=2em,itemsep=1pt]
\item every gridded price is a continuous martingale up to its expiry;
\item at every $T_l$ the spot sits on a node $K^{k_l}$ and every contract expiring at
  $T_l$ settles at its payoff;
\item if $k_l$ is interior, the model continues on $[T_l,T_{l+1}]$ from the finite state
  $z=\log q^{l+1}_{T_l}$ and $\lambda^{l+2},\dots,\lambda^M$ at $T_l$; if $k_l$ is an
  outer node, the spot is absorbed there: $q^j_t=e_{k_l}$ for all $j>l$ and $t\ge T_l$.
\end{enumerate}
\end{corollary}

\begin{proof}
By induction on $l$. On $[0,T_1]$ this is Proposition~\ref{prop:collapse} with
Theorem~\ref{thm:exist} and Remark~\ref{rem:horizon}, which also give the limit
$\lambda_{T_1}$, finite. Suppose the solution exists on $[0,T_l]$ with $\lambda_{T_l}$
finite and $q^l_{T_l}=e_k$.

If $k$ is interior, \eqref{eq:handover} and Lemma~\ref{lem:newfront} give a strictly
positive new front, so $z=\log q^{l+1}_{T_l}$ is finite, and
$K^{\!\top}q^{l+1}_{T_l}=K^k=S_{T_l}$, so the spot is continuous. On $[T_l,T_{l+1}]$ the
ladder is that of Sections~\ref{sec:drift}--\ref{sec:exist} with front $q^{l+1}$ and
expiries $T_{l+1}<\dots<T_M$, started from the $\mathcal F_{T_l}$-measurable state
$(z,\lambda^{l+2},\dots,\lambda^M)$. Theorem~\ref{thm:exist} applies with this
$\mathcal F_{T_l}$-measurable initial condition, the solution of
\cite[Theorem~V.38]{Protter2005} depending continuously on its starting point,
and with Proposition~\ref{prop:collapse} and Remark~\ref{rem:horizon} gives a unique
strong solution on $[T_l,T_{l+1}]$ that settles at $T_{l+1}$ with $\lambda_{T_{l+1}}$
finite. The densities $q^j$, $j>l+1$, are given by the same recursion on both sides of
$T_l$ with continuous inputs, so they are continuous at $T_l$.

If $k$ is an outer node, Lemma~\ref{lem:newfront} gives $q^j_{T_l}=e_k$ for every $j>l$.
The constant process $q^j_t\equiv e_k$ satisfies every recursion, since $A_je_k=e_k$, and
is a martingale, so all prices are constant and settle at their payoffs; we freeze
$\lambda$ at $\lambda_{T_l}$, on which no price then depends.

Uniqueness on each interval gives uniqueness on $[0,T_M]$.
\end{proof}

\section{Finite-step simulation and numerical evidence}\label{sec:numerics}

The drift restriction of Section~\ref{sec:drift} admits two numerically distinct
implementations.  The first integrates the log-DLV SDE directly.  It is the
closest implementation of the continuous-time model, but a finite time step is
only approximately martingale preserving.  The second advances the probability
density itself through a finite set of symmetric branches.  It is an exact
one-step martingale construction on the modeled finite support, although it is
not a claim about a continuous-strike or unbounded-support limit.  We report
both because their numerical guarantees differ.

All experiments in this section use four weekly expiries and fixed cash strikes
within each weekly transition.  ``DLVsq'' denotes the positive discrete local
variance $\dlv$.  The final density-branch audit uses the discrete factor $X$
alone.  Earlier development and long-horizon audits also multiply $X$ by an
independent mean-one Black factor $Y$ with annual volatility 10\%; the
conditioning identities below apply to both, using the Black semigroup for
$Y$.

\subsection{Direct integration of the DLV diffusion}\label{sec:numerics-sde}

For the direct SDE experiment we work in log excess-DLV coordinates
\[
  z^{j,i}=\log(\dlv^{j,i}-\dlv_{\min}).
\]
The volatility of $z$ is the rank-25 residual covariance estimated from 703
weekly historical transitions through 31 December 2024.  At every substep the
density volatility is obtained from \eqref{eq:volsolve}, and the drift of
$\dlv$ is then computed from \eqref{eq:pointwise}, including the It\^o term in
log coordinates.  Each week is integrated by ten stochastic-Heun substeps,
using the same Gaussian increment in predictor and corrector.  At the weekly
endpoint spot is drawn from the front density and the expired slab is removed.

The audit starts from historical state 44 (8 March 2022), uses 200 interior
strikes, 100,000 paths arranged as 50,000 antithetic pairs, and follows a
shrinking tower for four weeks without re-gridding.  Fifty antithetic pairs are
aggregated into each of 1,000 independent randomized-stratification blocks.
There are 800, 600, 400, and 200 live fixed-contract option returns at the four
trade dates, hence 2,000 coordinates in the joint test.

For this experiment the blockwise coordinate statistics are supplemented by a
10,000-draw shared-sign multiplier bootstrap.  If $t_k$ is the usual block
t-statistic, the two simultaneous statistics are
\[
  T_{\max}=\max_k|t_k|,
  \qquad
  T_{\mathrm{RMS}}=\left(d^{-1}\sum_{k=1}^{d}t_k^2\right)^{1/2}.
\]
The reported decision value is
$p_{\mathrm{comb}}=\min\{1,2\min(p_{\max},p_{\mathrm{RMS}})\}$.
The shared multiplier for all coordinates of a block preserves their observed
cross-option dependence.

\begin{table}[htbp]
\centering
\small
\begin{tabular}{lrrrrl}
\toprule
Audit & Active & $\max|t|$ & RMS $t$ & $p_{\mathrm{comb}}$ & 5\% decision\\
\midrule
One-step option returns       & 2,000 & 3.334 & 1.109 & 0.2746 & do not reject\\
Terminal option returns       & 2,000 & 3.334 & 1.044 & 0.3210 & do not reject\\
Terminal option price changes & 2,000 & 2.934 & 1.156 & 0.3500 & do not reject\\
Integrated-DLV control        & 1,200 & 1.950 & 0.878 & 1.0000 & do not reject\\
Weekly spot returns           & 4     & 2.431 & 1.561 & 0.0876 & do not reject\\
\bottomrule
\end{tabular}
\caption{Martingale diagnostics for the ten-substep Heun implementation.  The
800 omitted integrated-DLV coordinates have effectively zero block variance.
The primary one-step audit has separate multiplier-bootstrap values
$p_{\max}=0.1373$ and $p_{\mathrm{RMS}}=0.2465$; its simultaneous 5\% max-$t$
cutoff is 4.235.}
\label{tab:sde-martingale}
\end{table}

The largest density-mass residual was $9.33\times10^{-14}$ and the largest
conditional-forward residual was $8.93\times10^{-14}$.  The minimum density was
zero, due to finite-support point masses rather than a negative density.  The
simulation took 72.1 seconds, or approximately 1,388 paths per second.  These
results support the discretization, but they do not turn an infinitesimal drift
restriction into an exact one-week transition.

\subsection{An exact finite-step density construction}\label{sec:density-step}

We next construct the one-week law in density coordinates.  Let
$\varepsilon_b\in\{-1,+1\}^{d}$ be the $b$-th Rademacher branch and let
$\pi_b=2^{-d}$, with $d=4$ in the experiment.  Write $\Gamma^j$ for the
state-dependent loading of density $q^j$ implied by the chosen DLVsq loading.
The branch marginal is
\begin{equation}\label{eq:density-branch}
  q^{j,b}=q^j+\sqrt h\,\Gamma^j\varepsilon_b,
  \qquad b=1,\ldots,2^d .
\end{equation}
Since
$\sum_b\pi_b\varepsilon_b=0$ and
$\sum_b\pi_b\varepsilon_b\varepsilon_b^{\!\top}=I_d$,
\begin{equation}\label{eq:density-branch-mean}
  \sum_b\pi_bq^{j,b}=q^j
\end{equation}
holds exactly, before Monte Carlo sampling.  DLVsq is then recovered
algebraically from each $q^{j,b}$ by Proposition~\ref{prop:codebook}; it is not
advanced by an Euler step or fitted by a nonlinear optimizer.  A branch is
accepted only if its masses are non-negative and the recovered DLVsq lies in
its configured open bounds.  The implementation checks mass, mean, decoder
consistency, and DLV bounds at every step.

In the implemented weekly law the front density is branch invariant,
$q^{1,b}=q^1$.  Independently of $b$, a front node $I=i$ is drawn with
probability $q^1_i$.  For every surviving expiry $j\ge2$, let
$R^{b,j}_{ir}$ be the conditional probability of node $r$ given branch $b$ and
front node $i$.  The SANOS transition supplies rows satisfying
\begin{align}
  \sum_rR^{b,j}_{ir}&=1,
  &\sum_rx_rR^{b,j}_{ir}&=x_i,
  &\sum_iq^1_iR^{b,j}_{ir}&=q^{j,b}_r .
  \label{eq:conditional-density-identities}
\end{align}
The first two identities make every row a martingale transition; the third
recovers the branch marginal.  Together with
\eqref{eq:density-branch-mean}, they imply the tower property for every payoff
represented on the finite support.

This construction is exact on the finite model grid.  It is not asserted to be
the unique risk-neutral law, and the particular independent sampling of $b$
and $I$ does not generate contemporaneous correlation between the raw surface
branch and front-node innovation.  Realized density-weighted factor covariance
nevertheless matched the intended empirical covariance closely: in the
independent all-state confirmation the largest relative error was 1.87\% and
the RMS error was 0.55\%.

\subsection{Conditioning and re-gridding}\label{sec:numerics-regrid}

Suppose current cash spot is $s$ and the revealed front node is $x_i$, so the
next $X$ spot is $s x_i$.  Define the next relative support by
\begin{equation}\label{eq:exact-regrid}
  x_r^{\prime(i)}=\frac{x_r}{x_i}.
\end{equation}
Then $(s x_i)x_r^{\prime(i)}=s x_r$ path by path.  Consequently the price of a
surviving fixed-cash call is identical on the inherited and re-gridded
coordinates,
\[
  C_{t+h}^{b,i,j}(K)
  =\sum_rR^{b,j}_{ir}(s x_r-K)^+.
\]
Averaging first over $i$ and then over $b$ gives
\begin{align*}
  \E_t[C_{t+h}^{j}(K)]
  &=\sum_b\pi_b\sum_iq^1_i\sum_rR^{b,j}_{ir}(s x_r-K)^+\\
  &=\sum_rq^j_r(sx_r-K)^+=C_t^j(K).
\end{align*}
Thus re-gridding is an exact coordinate change for all surviving fixed-cash
claims, not an interpolation.  With an independent Black factor $Y$, the same
proof conditions on $Y$ and then applies its mean-one Black semigroup.  A newly
issued far expiry is different: there is no pre-existing fifth-week claim whose
price it must preserve.

At each weekly roll, the old 2W--4W densities become the new floating 1W--3W
densities through \eqref{eq:conditional-density-identities} and
\eqref{eq:exact-regrid}.  Only the new 4W DLVsq curve is supplied statistically.
The retained boundary model is a 25-factor PCA followed by an unrestricted
weekly VAR(1) in factor scores.  It uses 703 training pairs through 2024,
explains 99.9912\% of the 4W curve variation, and has spectral radius 0.9089.
On 155 untouched 2025 pairs its raw-coordinate RMSE is 0.5902, against 0.6526
for persistence, a 9.56\% improvement.  SANOS constructs the new 4W density as
a static martingale extension of the inherited 3W density.

\subsection{Martingale audit for the exact density step}\label{sec:density-tests}

Adjacent simulated paths form an antithetic pair.  For pair $a=1,\ldots,n$ and
audited return coordinate $k$, let $R_{ak}$ be the pair mean, $\bar R$ the
vector of sample means, and $\widehat\Sigma$ their sample covariance.  We
report
\begin{equation}\label{eq:density-wald}
  z_k=\frac{\sqrt n\,\bar R_k}{\widehat\sigma_k},
  \qquad
  J=n\bar R^{\!\top}\widehat\Sigma^{+}\bar R .
\end{equation}
The pseudo-inverse is formed in correlation space after discarding eigenmodes
at or below 0.01.  If $\nu$ modes remain, the reference value is
$J\sim\chi^2_\nu$ under the joint zero-mean null.  A small upper-tail $p$-value
rejects zero drift; a large value means only that the observed drift is
compatible with Monte Carlo noise.  We use 0.05 for one look, 0.0125 for four
weekly looks, and 0.003125 for 16 looks.  Because the options are highly
collinear and the retained modes are estimated from the sample, the exact
density-tower identities and matched cash/re-grid prices are stronger
structural diagnostics than an isolated asymptotic $p$-value.

The canonical X-only audit used all 891 certified historical states, 100,000
paths (50,000 independent antithetic pairs), 11 ITM call/put strikes per
expiry, four one-week steps, 64-bit arithmetic, and 10\% annualized log-DLVsq
volatility.  Each state received 56 or 57 independent pairs.  The audited book
contains spot and the surviving original fixed-cash claims.

\begin{table}[htbp]
\centering
\small
\begin{tabular}{rrrrrr}
\toprule
Week & Worst $|z|$ & Max option mean (bp) & $J$ & $\nu$ & $p$\\
\midrule
1 & 1.796 & 0.497 & 10.980 & 7 & 0.139\\
2 & 2.143 & 0.863 &  7.228 & 7 & 0.406\\
3 & 1.371 & 0.696 &  5.858 & 6 & 0.439\\
4 & 2.057 & 1.245 &  5.506 & 5 & 0.357\\
\bottomrule
\end{tabular}
\caption{All-state X-only martingale audit for the surviving original
fixed-cash book plus spot.  Option means are basis points of initial spot.  No
horizon rejects the joint zero-mean null at 5\%.}
\label{tab:density-all-state}
\end{table}

The deterministic checks close much more tightly than Monte Carlo sampling:
the maximum density-tower residual is $6.49\times10^{-15}$, option-tower
residual $1.67\times10^{-16}$, spot-mean residual $2.60\times10^{-14}$,
four-roll decoder residual $1.77\times10^{-13}$, and first re-grid cash-price
residual $3.33\times10^{-16}$.

A matched run decomposed every weekly change into the fixed-expiry surface
branch and the spot-conditioning/re-grid component, using identical random
draws.  Direct inherited-cash and re-gridded valuations agree pathwise to at
most $2.22\times10^{-16}$ and have identical $z/J/p$ statistics.

\begin{table}[htbp]
\centering
\small
\begin{tabular}{rrrrr}
\toprule
Week & Fixed-expiry $p$ & Full first-member $p$ & Full pair $p$ & Max price difference\\
\midrule
1 & 0.366 & 0.684  & 0.139  & 0\\
2 & 0.538 & 0.163  & 0.128  & $2.22\times10^{-16}$\\
3 & 0.911 & 0.414  & 0.0208 & $2.22\times10^{-16}$\\
4 & 0.945 & 0.0808 & 0.546  & $2.22\times10^{-16}$\\
\bottomrule
\end{tabular}
\caption{Matched fixed-support and floating-coordinate decomposition.  The
first two $p$-value columns use one independent member of each pair; the third
uses antithetic pair means.}
\label{tab:density-regrid}
\end{table}

The week-3 pair-mean value is a nominal rejection at 5\%, but not at the
pre-specified four-look cutoff 0.0125.  It is also visibly sensitive to the
correlation eigenvalue cutoff: its $p$-value is 0.0429, 0.0208, 0.0879, and
0.295 at cutoffs 0.005, 0.01, 0.02, and 0.05.  Estimating covariance and the
mean on opposite half-samples gives 0.214 and 0.338.  No coordinate has
$|z|>2$ and the largest mean is 0.363bp.  Since the direct and re-gridded path
matrices are identical, this fluctuation cannot be caused by re-gridding.
The fixed-expiry branch alone has maximum option means of only 0.0046, 0.0056,
0.0013, and 0.0050bp over weeks 1--4.

Earlier development audits reached the same conclusion under a nonzero Black
backbone.  On 21 states with 10,000 paths per state, the four weekly maximum
coordinate $|z|$ values were 2.400, 3.203, 2.871, and 2.968, while the
covariance-aware Fisher-combined state-test $p$-values were 0.871, 0.672,
0.852, and 0.522.  A heterogeneous 10,000-path confirmation covering all 891
states gave unconditional maximum $|z|$ values 1.575, 1.970, 1.878, and 1.040
and $p$-values 0.702, 0.414, 0.067, and 0.317.  Its high-dimensional
state-feature-weighted values were 0.027, 0.759, 0.251, and 0.0089; they are
reported because they expose the poor calibration of hard spectral cutoffs in
small per-state samples, rather than being hidden behind the aggregate result.

\subsection{Repeated rolls, speed, and finite-support behavior}

A 16-week audit with the independent 10\% Black backbone used 10,000 paths and
all 891 states.  Across all rolls the largest decoder residual was
$2.35\times10^{-13}$, conditional mass and mean errors stayed below
$6.6\times10^{-14}$, and DLVsq remained strictly inside $(10^{-4},9)$.  The
maximum spot $|z|$ over 16 horizons was 2.126, the maximum rolling-option
$|z|$ was 2.417, and every rolling-option joint $p$-value was at least 0.0839.
With the last-expiry AR boundary, median cross-expiry dispersion of
$\sqrt{\mathrm{DLVsq}}$ was 0.0459 after week 1 and 0.0569 after week 16, so
the term structure did not collapse.  The corresponding implementation ran
160,000 path-week steps in 14.6 seconds.

The exact branch method is materially faster than a nonlinear finite-step
solve.  In the heterogeneous 10,000-path benchmark it processed 40,000
path-week steps in 8.84 seconds (4,527 per second), including 7.36 seconds to
construct 5,000 distinct initial kernels.  This was about 23 times faster than
the earlier finite-$h$ Newton transition (195 per second), but about 5.7 times
slower than the local SDE transition (25,649 per second).  The final X-only
audit, with more extensive all-state diagnostics, processed 400,000 path-week
steps in 169.8 seconds.  These timings are indicative rather than a controlled
hardware benchmark.

The remaining limitation is economic rather than a floating-point failure.  If
the lowest support node $x_0$ is revealed, then
\[
  \sum_rR^{b,j}_{0r}x_r=x_0,
  \qquad x_r\ge x_0,
\]
forces $R^{b,j}_{00}=1$.  The upper endpoint is analogous.  The $X$ law is then
absorbing, all surviving options are intrinsic, and interior DLVsq coordinates
are price-unidentified.  The positive-density formula $\eta=\psi/q$ no longer
applies at zero-mass nodes; any continued DLV coordinates there are an
economically invisible numerical convention.

In the four-week X-only run, boundary incidence was 0.003\%, 0.003\%, 0.003\%,
and 0.006\%.  In the 16-week run, the fully collapsed fraction grew from
0.01\% to 0.48\%; the median week-16 front density retained 23.44 effective
nodes, while its first percentile retained 1.15.  Avoiding this absorbing tail
requires dynamic support expansion, an explicit tail continuation, or a
boundary state carrying a conditional tail.  Re-gridding alone cannot create
support beyond an extreme node while preserving the inherited option surface.

The numerical evidence therefore supports three distinct statements: the
Heun implementation is a credible approximation to the continuous density
drift; the symmetric density construction is exactly martingale consistent on
the finite support; and weekly re-gridding preserves every surviving
fixed-cash price path by path.  It does not establish uniqueness, an
unbounded-support model, or the absence of boundary bias at horizons much
longer than those tested.

\section{What is open}\label{sec:open}

\begin{enumerate}[leftmargin=1.6em,itemsep=0.4em]
\item \emph{An off-lattice terminal spot.} By Remark~\ref{rem:bary} the chart
represents every $s\in[K^0,K^{N+1}]$ at expiry, and the state that does so is
$\varpi(s)$, carried by the two nodes straddling $s$. Whether a front volatility
exists driving $q^1$ to $\varpi(s)$ for $s$ strictly inside a cell, rather than to
a single node as \eqref{eq:nufront} does, is open. The obstruction is that
exact settlement forces $q^1_{T_1}=\varpi(S_{T_1})$ and hence
$\int_0^{T_1}(\eta^{1,i})^2\dd t<\infty$ for both surviving nodes, which by the
argument of Proposition~\ref{prop:collapse} they cannot both satisfy when the
volatility is proportional to distance from the spot. Until it is answered the
model of Proposition~\ref{prop:collapse} pins, as \cite{Wissel2007} does, and
his question of a front that does not pin the spot remains open.

\item \emph{The continuum limit.} The constant $(M-1)(N+1)$ grows with the grid. Whether $\sup_N\Theta_j<\infty$ on a class of
realistic surfaces, which is what a continuum statement would need, is open.
\end{enumerate}


\bibliographystyle{plain}
\bibliography{refs}

\appendix
\section{Relation to the continuum and lattice results}\label{app:priority}

\emph{TBD.}

\subsection{Carmona and Nadtochiy}\label{app:cn}

\emph{TBD.}

\subsection{Wissel}\label{app:wissel}

\emph{TBD.}

\begin{remark}[The quantile chart forces the spot onto the grid]\label{rem:pinning}
\emph{TBD.}
\end{remark}


\end{document}
